% Ad Atticum 1.8 (ad-atticum-01-08)
% Source: https://raw.githubusercontent.com/PerseusDL/canonical-latinLit/master/data/phi0474/phi057/phi0474.phi057.perseus-lat1.xml
% Fetched by scripts/fetch_latin.py

% Dateline (Perseus): Scr. Romae post Id. Febr. a. 687 (67) .

\ciceroLetterOpener{CICERO ATTICO salutem}

\ciceroSection{1} apud te est ut volumus. mater tua et soror a me Quintoque fratre diligitur. cum Acutilio sum locutus. is sibi negat a suo procuratore quicquam scriptum esse et miratur istam controversiam fuisse quod ille recusarit satis dare amplius abs te non peti. quod te de Tadiano negotio decidisse scribis, id ego Tadio et gratum esse intellexi et magno opere iucundum. ille noster amicus, vir me hercule optimus et mihi amicissimus, sane tibi iratus est. hoc si quanti tu aestimes sciam, tum quid mihi elaborandum sit scire possim.

\ciceroSection{2} L. Cincio HS CCIↃↃ CCIↃↃ CCCC pro signis Megaricis, ut tu ad me scripseras, curavi. Hermae tui Pentelici cum capitibus aeneis, de quibus ad me scripsisti iam nunc me admodum delectant. qua re velim et eos et signa et cetera quae tibi eius loci et nostri studi et tuae elegantiae esse videbuntur quam plurima quam primumque mittas et maxime quae tibi gymnasi xystique videbuntur esse. nam in eo genere sic studio efferimur ut abs te adiuvandi, ab aliis prope reprehendendi simus. si Lentuli navis non erit, quo tibi placebit imponito.

\ciceroSection{3} Tulliola deliciolae nostrae tuum munusculum flagitat et me ut sponsorem appellat; mi autem abiurare certius est quam dependere.
